\documentclass[11pt,a4paper]{article}
\usepackage[utf8]{inputenc}
\usepackage[T1]{fontenc}
\usepackage[spanish]{babel}
\usepackage{graphicx}
\usepackage{float}
\usepackage[margin=2.2cm]{geometry}
\usepackage{enumitem}
\usepackage{booktabs}
\usepackage[hidelinks]{hyperref}

\title{Guía de observación\\EDA 03: cobertura Landsat y rendimiento 2025}
\author{Proyecto AgroCebada}
\date{}

\newcommand{\edafigure}[3]{
  \subsection{#1}
  \begin{figure}[H]
    \centering
    \includegraphics[width=0.96\textwidth,height=0.66\textheight,keepaspectratio]{\detokenize{#2}}
    \caption{#1}
  \end{figure}
  \textbf{Preguntas guía}
  #3
}

\begin{document}
\maketitle
\tableofcontents
\newpage

\section{Propósito y definición de cobertura}
Este EDA estudia por qué Landsat tiene cobertura desigual entre parcelas y cómo
representar esa desigualdad antes del modelado. La cobertura histórica se calcula
respecto a las 392 fechas Landsat presentes entre 2022 y 2025. También se calcula
la cobertura específica del ciclo objetivo, abril--octubre de 2025, que contiene
53 fechas disponibles.

El criterio principal de agrupación es la proporción de fechas históricas con
\texttt{vi6t\_promedio} válido. Los grupos se forman con los cuartiles de esta
proporción: baja hasta 0.3597, media entre 0.3597 y 0.4158, y alta por encima de
0.4158. Es una segmentación descriptiva y reproducible, no un resultado de un
modelo de clustering.

\section{Cobertura por parcela}

\edafigure{Cobertura Landsat utilizable ordenada}
{../../../outputs/eda/01_tablas_maestras/03_cobertura_landsat/01_cobertura_relativa_ordenada.png}{
\begin{itemize}[nosep]
  \item ¿La cobertura cambia gradualmente o existen grupos separados?
  \item ¿Cuántas parcelas pertenecen a cada nivel?
  \item ¿Los umbrales representan calidad agronómica o sólo disponibilidad relativa?
\end{itemize}}

\edafigure{Distribución de cobertura total y utilizable}
{../../../outputs/eda/01_tablas_maestras/03_cobertura_landsat/02_distribucion_cobertura.png}{
\begin{itemize}[nosep]
  \item ¿Cuánto disminuye la cobertura al exigir un VI6T válido?
  \item ¿Hay parcelas con muchas capturas pero pocos valores utilizables?
  \item ¿Conviene usar número de filas o número de valores válidos como indicador de confianza?
\end{itemize}}

\edafigure{Cobertura histórica frente al ciclo 2025}
{../../../outputs/eda/01_tablas_maestras/03_cobertura_landsat/03_historica_vs_ciclo_2025.png}{
\begin{itemize}[nosep]
  \item ¿La cobertura histórica predice la disponibilidad durante abril--octubre de 2025?
  \item ¿Existen parcelas con buena historia pero mala cobertura en el ciclo objetivo?
  \item ¿Qué cobertura debe tener prioridad para predecir el rendimiento 2025?
\end{itemize}}

\edafigure{Cobertura mensual de Landsat}
{../../../outputs/eda/01_tablas_maestras/03_cobertura_landsat/04_heatmap_cobertura_mensual.png}{
\begin{itemize}[nosep]
  \item ¿La baja cobertura se concentra en meses o parcelas específicas?
  \item ¿Hay bloques persistentes que indiquen grupos de adquisición distintos?
  \item ¿Las parcelas con menos datos conservan observaciones dentro del ciclo 2025?
\end{itemize}}

\edafigure{Huecos temporales por grupo}
{../../../outputs/eda/01_tablas_maestras/03_cobertura_landsat/05_huecos_temporales_por_grupo.png}{
\begin{itemize}[nosep]
  \item ¿Los grupos con menor cobertura también tienen huecos máximos mayores?
  \item ¿La diferencia principal es continuidad o cantidad total de observaciones?
  \item ¿Los huecos son suficientemente grandes para impedir agregaciones mensuales?
\end{itemize}}

\section{Relación con atributos de parcela y objetivo}

\edafigure{Cobertura utilizable frente al área}
{../../../outputs/eda/01_tablas_maestras/03_cobertura_landsat/06_cobertura_vs_area.png}{
\begin{itemize}[nosep]
  \item ¿Las parcelas pequeñas reciben sistemáticamente menor cobertura?
  \item ¿El área explica los grupos observados?
  \item ¿Debe el área formar parte del control de calidad o del modelo agronómico?
\end{itemize}}

\edafigure{Cobertura utilizable frente al rendimiento}
{../../../outputs/eda/01_tablas_maestras/03_cobertura_landsat/07_cobertura_vs_rendimiento.png}{
\begin{itemize}[nosep]
  \item ¿Existe asociación entre cobertura y rendimiento en entrenamiento?
  \item ¿La asociación puede deberse a geografía, manejo, fenología o calidad de medición?
  \item ¿Sería seguro usar la cobertura como peso de entrenamiento?
\end{itemize}}

\edafigure{Distribución espacial de los grupos}
{../../../outputs/eda/01_tablas_maestras/03_cobertura_landsat/08_mapa_grupos_cobertura.png}{
\begin{itemize}[nosep]
  \item ¿Los grupos de cobertura forman conglomerados espaciales?
  \item ¿La cobertura podría estar actuando como sustituto de municipio, órbita o región?
  \item ¿Cómo debe diseñarse la validación para evitar aprovechar accidentalmente esa estructura espacial?
\end{itemize}}

\section{Respuestas y hallazgos}

\subsection{Grupos y disponibilidad}
\begin{enumerate}
  \item La distribución no es completamente continua: aparece un grupo superior
  con cobertura claramente mayor. Hay 51 parcelas en cobertura baja, 97 en media
  y 49 en alta. Los límites son cuartiles del conjunto actual y no umbrales de
  calidad agronómica universal.

  \item La mediana de observaciones históricas es 205, 208 y 296 para los grupos
  bajo, medio y alto. Sin embargo, las medianas de VI6T válido son solamente 136,
  148 y 184. Por ello, contar filas sobreestima la información utilizable. La
  proporción de VI6T válido respecto a las 392 fechas disponibles es el indicador
  de confianza más directo de los evaluados.

  \item En abril--octubre de 2025, la cobertura mediana relativa es 0.491 en los
  grupos bajo y medio y 0.755 en el alto. La historia ayuda a anticipar la
  disponibilidad del ciclo objetivo, pero no la sustituye. Para predecir el
  rendimiento 2025 debe priorizarse la cobertura dentro de abril--octubre; la
  cobertura 2022--2024 debe considerarse contexto histórico.

  \item El heatmap indica que la diferencia es persistente por parcela y no la
  ausencia de un único mes global. Incluso los grupos bajos conservan datos en el
  ciclo objetivo, de modo que pueden producir agregados mensuales, aunque con
  menor estabilidad y diferente número de observaciones.

  \item Los huecos máximos medianos son muy parecidos: 16 días en los grupos bajo
  y medio y 15 en el alto. La separación entre grupos se debe principalmente al
  número total de capturas y valores VI6T válidos, no a interrupciones extremas
  mucho mayores. Huecos de este tamaño permiten agregación mensual, pero deben
  registrarse el número de observaciones y los días cubiertos.
\end{enumerate}

\subsection{Área, rendimiento y espacio}
\begin{enumerate}[resume]
  \item El área casi no explica la cobertura utilizable: la correlación de
  Spearman es $-0.027$. Las medianas de área son 4.13 ha en cobertura baja,
  6.89 ha en media y 3.79 ha en alta, sin un gradiente consistente. El área puede
  ser una variable agronómica, pero no debe utilizarse como sustituto de calidad.

  \item Sí existe una asociación descriptiva entre cobertura y rendimiento. La
  correlación Spearman es 0.290 para cobertura VI6T válida y 0.378 para número de
  observaciones. El rendimiento medio por grupo es 3.72, 3.76 y 4.55 t/ha; las
  medianas son 3.50, 3.50 y 4.54. Esto no significa que observar más eleve el
  rendimiento: la cobertura puede estar capturando región, órbita, condiciones
  del cultivo o diferencias del proceso de datos.

  \item Los grupos contienen parcelas de ambos conjuntos: baja tiene 33 de
  entrenamiento y 18 de predicción; media, 69 y 28; alta, 36 y 13. No hay un
  grupo exclusivo del conjunto oculto. El mapa debe usarse para determinar si el
  grupo alto se concentra espacialmente; si es así, una partición aleatoria puede
  producir resultados demasiado optimistas y conviene validar por grupos
  espaciales o municipios.
\end{enumerate}

\section{¿Debe darse peso a las parcelas según su cobertura?}
Es razonable probarlo, pero no asumir que es correcto. El archivo de métricas
incluye un peso candidato de cobertura, definido como la proporción histórica
de VI6T válido dividida entre la mayor proporción observada. Este peso expresa
disponibilidad, no certeza del rendimiento ni importancia agronómica.

Se recomiendan tres experimentos bajo las mismas particiones de validación:
\begin{enumerate}
  \item modelo sin peso, incorporando cobertura y número de observaciones como variables;
  \item modelo con peso candidato durante el entrenamiento;
  \item agregados temporales acompañados por su conteo, dispersión y cobertura, sin ponderar la pérdida.
\end{enumerate}
Debe elegirse la opción por desempeño fuera de muestra en cada grupo de cobertura,
no sólo por el error global. Ponderar hacia abajo las parcelas de baja cobertura
podría perjudicar precisamente a las 18 parcelas de predicción de ese grupo.

\section{Qué representa el rendimiento}
La documentación oficial indica que \texttt{RENDIMIENTO\_T\_HA} es el rendimiento
de cebada en toneladas por hectárea correspondiente al ciclo productivo
abril--octubre de 2025. La tabla contiene una fila por parcela, no una serie
temporal: 138 parcelas de \texttt{ENTRENAMIENTO} tienen rendimiento y las 59 de
\texttt{PREDICCION} lo mantienen oculto.

No se proporciona una fecha exacta de cosecha por parcela. ``Abril--octubre de
2025'' es una ventana de ciclo, no evidencia de que todas las parcelas hayan sido
cosechadas el mismo día. Tampoco hay rendimientos etiquetados para 2022, 2023 o
2024. Esos años pueden aportar climatología o comportamiento histórico, pero no
constituyen ejemplos adicionales de entrenamiento supervisado.

Para evitar fuga de información, la construcción de variables debe fijar un
momento de predicción. Si el objetivo es estimar al final del ciclo, pueden usarse
capturas hasta octubre de 2025. Si se desea predicción anticipada, deben definirse
cortes reproducibles, por ejemplo hasta junio, julio, agosto o septiembre, y
eliminar toda observación posterior a cada corte.

\section{Conclusión operativa}
Sentinel-2 puede ser la columna vertebral por uniformidad y riqueza de variables.
Landsat sigue siendo útil como señal complementaria: su cobertura no invalida el
sensor, pero exige conservar métricas de disponibilidad, agregar dentro de
ventanas comunes y evaluar desempeño por grupo. El peso por cobertura debe ser
un experimento de validación, no una corrección automática.

\end{document}
